\section{晋：中原霸权与卿族的影子}

晋在春秋中期最像一个能真正组织北方秩序的大国：腹地较大，卿族能带兵，国君能借流亡、婚姻和会盟建立关系。但它的强大也不只属于国君。军队、封邑和政治资源被多支卿族掌握，国家每扩张一步，内部分享权力的人也更多一步。

\begin{event}{公元前636年}{重耳返晋，即位为文公}{可靠纪年}{晋}\eventanchor{SA-0636-JIN-WEN}{春秋·晋文公即位}
\term{这一年发生了什么}：重耳结束长期流亡返晋，即位为晋文公。

\term{事情为什么重要}：流亡不是一段只为人物增色的传奇。重耳在齐、楚、秦等国长期活动，形成了跨国的人际和政治网络；这些关系会在\eventref{SA-0632-CHENGPU}{春秋·城濮之战}时转化为联盟资源。

\term{史料边界}：《左传》和《史记·晋世家》都把流亡写得极富人物性格。它们可以说明当时人如何理解政治信誉与结盟，不应被当成每一次对话的逐字记录。
\end{event}

\begin{event}{公元前632年}{城濮之战，晋文公成为北方霸主}{可靠纪年}{城濮}\eventref{SA-0632-CHENGPU}{春秋·城濮之战}
\term{这一年发生了什么}：晋与齐、宋、秦等联盟在城濮击败楚军。

\term{后来改变了什么}：晋能够把多国关系组织成可行动的联盟，成为中原会盟的关键。霸权由此不是单纯的“晋打赢楚”，而是晋文公的外交网络、北方资源和楚方联盟脆弱性共同产生的结果。
\end{event}

\begin{event}{公元前597年}{邲之战，晋军败于楚}{可靠纪年}{邲}\eventref{SA-0597-BI}{春秋·邲之战}
\term{这一年发生了什么}：晋楚在邲交战，晋军败退。

\term{事情为什么重要}：城濮后的晋并没有永久压住楚。战场上暴露的指挥分歧，也折射出晋国内部卿大夫对军政资源的不同控制。楚的胜利是晋楚两极重新平衡，而不是晋霸从此消失。
\end{event}

\begin{event}{公元前575年}{鄢陵之战}{可靠纪年}{鄢陵}\eventanchor{SA-0575-YANLING}{春秋·鄢陵之战}
\term{这一年发生了什么}：晋楚再次交战，晋取得胜利。

\term{后来改变了什么}：鄢陵表明晋仍有军事组织能力，但连续战争并没有让任何一方获得决定性解决。大国在长期消耗后，转向\eventref{SA-0546-MIBING}{春秋·弭兵之会}式的均势安排。
\end{event}

\begin{event}{公元前546年}{弭兵之后，卿族问题走向中心}{可靠纪年}{晋及中原}\eventref{SA-0546-MIBING}{春秋·弭兵之会}
\term{这一阶段发生了什么}：对外大规模战争减少，晋国内部的卿族竞争更突出。

\term{后来改变了什么}：外部压力暂缓并未让晋国更稳定，反而让掌握封邑、军队和家臣网络的卿族有更长时间经营自身。最终的\eventref{WQ-0453-ZHI}{战国·三家灭智}不是突发政变，而是这条内部权力线的终点。
\end{event}

\begin{sourcenote}[晋国为何会“强而后分”]
《左传》对晋国卿大夫政治有大量记述，但它往往以人物伦理和礼法冲突组织故事。把这些故事与《史记·晋世家》、后来的三家分晋年表相互对读，可以看见更长的结构：国君依靠卿族扩张，卿族也在扩张中获得独立资源。不要把晋的分裂简化成某一位国君失德或某一家族野心。
\end{sourcenote}

\begin{center}
\begin{tikzpicture}[x=.88cm,y=.88cm,font=\small]
  \fill[paper] (0,0) rectangle (10.5,5.7);
  \draw[blue!55,line width=1pt] (.2,4.8) .. controls (3.2,4.6) and (6.3,5.1) .. (10.0,4.6);
  \node[text=blue!70!black,above] at (5.0,4.9) {黄河};
  \draw[blue!45,line width=.8pt] (3.8,.4) .. controls (5.0,2.0) and (6.2,2.2) .. (8.8,.7);
  \node[text=blue!70!black,below] at (6.2,1.35) {淮水方向};

  \fill[green!20] (1.0,3.5) -- (4.0,3.3) -- (4.2,5.0) -- (1.6,5.2) -- cycle;
  \node[align=center] at (2.5,4.25) {晋\\北方联盟};
  \fill[purple!18] (4.8,.8) -- (8.4,.9) -- (8.7,3.4) -- (5.5,3.6) -- cycle;
  \node[align=center] at (6.7,2.2) {楚\\北上军队};
  \fill[yellow!25] (7.6,3.5) -- (9.4,3.4) -- (9.6,4.7) -- (7.9,4.9) -- cycle;
  \node at (8.6,4.1) {齐};
  \fill[orange!24] (4.1,3.2) -- (5.1,3.2) -- (5.2,4.1) -- (4.2,4.2) -- cycle;
  \node at (4.65,3.65) {宋};
  \fill[timeline] (5.1,3.0) circle (.14);
  \node[below,align=center] at (5.1,2.9) {城濮\\战场};

  \draw[-{Stealth[length=2mm]},ultra thick,color=green!55!black] (2.8,3.75) -- (4.95,3.05);
  \draw[-{Stealth[length=2mm]},ultra thick,color=dynasty] (6.1,2.6) -- (5.25,2.95);
  \node[text=green!55!black,above] at (3.75,3.35) {晋及盟军};
  \node[text=dynasty,below] at (5.65,2.55) {楚军北上};
\end{tikzpicture}

\smallskip
\small\textit{图  城濮战役空间示意：突出晋的北方联盟、楚的北上方向、宋等中原国家和战场的相对位置；不表示精确行军路线或疆界。}
\end{center}


\begin{turningpoint}[制度得失：分权军政的上升与终点]
晋的卿族分权让国家能同时动员多支军事和地方资源，这是其霸权的优势；代价是国君难以把胜利积累为独占的中央能力。春秋晋国的成功，恰好给战国韩、赵、魏的独立提供了人、地、兵和制度遗产。
\end{turningpoint}
